\section{Introduction}
\label{sec:introduction}

\subsection{Background and motivation}
\label{sec:motivation}

Fixed-wing unmanned aerial vehicle (UAV) mission planning involves coupled
decisions about route geometry, travel time and energy expenditure. In a
threat environment, these decisions also determine the duration and
intensity of exposure. Two collision-free routes can therefore have markedly
different mission costs: a detour may avoid a region of high threat intensity
but require more time and fuel, whereas a shorter route may cross that
region for a limited duration. A useful planner must distinguish absolute
prohibitions from traversable regions whose desirability depends on the
mission preference.

A* provides a general mechanism for searching a graph with assigned edge
costs~\cite{hart1968}. Its capacity to represent soft threats depends on the
chosen graph and cost model. In a binary-occupancy, time-oriented formulation,
obstacle checks establish admissibility and travel-time estimates rank the
remaining routes; different threat intensities within free space do not
affect that ranking unless they are included in the cost. The question
addressed here is how to combine this spatial information with time and
fuel preferences while retaining a clear connection to the trajectory
that will ultimately be evaluated.

Three issues arise in this connection. First, exposure is accumulated over
time, so a spatial threat score alone cannot determine the cost of a
variable-speed trajectory. Second, increasing the resolution of a
three-dimensional search expands its computational requirements, motivating
local refinement of a coarse route. Third, collision-free grid edges can
meet at corners that a fixed-wing aircraft cannot follow directly.
Smoothing changes the geometry, and speed assignment changes the time spent
along it. Consequently, the mission metrics of the final trajectory must
be distinguished from the estimates used to select the discrete path.

\subsection{Related work}
\label{sec:related_work}

Risk-aware and multi-objective planning incorporate more information than
binary occupancy. Guo et al.~\cite{guo2022} combine threat strength and path
length in FC-RRT*, using the resulting cost and flight constraints in tree
expansion and parent selection. Their formulation illustrates that
traversable threats and mission-dependent trade-offs are established
planning concepts. For the present problem, the modeling emphasis is on
time-integrated exposure together with flight time and fuel consumption.
Because subsequent speed assignment can alter exposure even when the route
is unchanged, these quantities also require a common final evaluation
protocol.

Hierarchical planning addresses the cost of resolving large spatial
domains. Liu et al.~\cite{liu2026} couple a multigranularity semantic-risk
representation with coarse corridor guidance and fine-level A* search.
Their urban framework also supports replanning after airspace changes; its
voxel-center output requires subsequent smoothing and platform-specific
dynamic verification. This work establishes a close precedent for
risk-aware corridor refinement. The present study concentrates on
fixed-wing trajectory generation after route selection, carrying shared
time--exposure--fuel references through both search resolutions and
evaluating the refined trajectory under its assigned speed profile.

Fixed-wing trajectory methods explicitly address the relation between
geometry and motion. Li et al.~\cite{li2025} decompose constrained trajectory
optimization into path and speed subproblems and use nonlinearity transfer
and convexification to obtain convergent solution procedures. Zhang et
al.~\cite{zhang2026} use an improved A* path to initialize states and controls
for subsequent pseudospectral trajectory optimization. These studies show
that path--speed decomposition and search-to-trajectory integration already
have substantive precedents. Here, a reduced flight-envelope formulation
is used to connect preference-based graph search with spline refinement
and speed assignment, rather than to solve a full aircraft optimal-control
problem. Parametric spline fitting itself follows established
methods~\cite{dierckx1993}; the application-specific issue is how the fitted
curve is constrained, selected and reevaluated.

\subsection{Research focus and contributions}
\label{sec:contributions}

The problem addressed in this paper is the consistency of the transition
from mission-dependent route selection to a constrained fixed-wing
trajectory. This consistency has three concrete requirements: the two
search resolutions must use the same objective references; the smoothed
curve must be checked against the constraints affected by geometric
modification; and final time, exposure and fuel must be calculated using
the assigned motion rather than inherited from the search. The proposed
framework addresses these requirements through the following contributions.

\begin{enumerate}
\item A mission-preference-driven time--exposure--fuel formulation is
constructed for hard--soft threat environments. Hard no-fly regions define
infeasibility, while a continuous radar-exposure surrogate enters the cost
through its time integral. Three single-objective coarse-search anchors
provide fixed reference scales for weighted route selection and subsequent
reporting.

\item A coarse-to-fine three-dimensional planning procedure applies the
same weighted search objective and reference scales at both resolutions.
The coarse route defines a spatial corridor within which a fresh fine-grid
search resolves local geometry. This procedure makes the role of spatial
refinement explicit without interpreting interpolation as replanning.

\item A constrained path-to-trajectory procedure connects cubic B-spline
refinement, analytical curvature evaluation, sampled feasibility checks
and forward--backward speed assignment. Final mission costs are
reintegrated along the refined trajectory, exposing the changes introduced
by geometric modification and retiming instead of treating discrete-path
statistics as executable-trajectory performance.
\end{enumerate}

The numerical study is organized around controlled comparisons in one
prescribed environment. Full two-stage cases examine the objective choice
and the successive planning stages, while a separate 72-case coarse-search
protocol characterizes preference sensitivity and supplies representative
trajectories. This design separates evidence about the complete planning
chain from evidence about weight-dependent behavior. Feasibility throughout
this paper refers to the specified reduced flight envelope and numerical
checks; it does not imply a full flight-dynamics or tracking validation.

Section~\ref{pr:sec:problem} defines the environment, constraints and mission
objective. Section~\ref{sec:method} presents the planning and evaluation
procedure. Section~\ref{sec:results} specifies the controlled comparisons
and preference analysis, and Section~\ref{sec:conclusions} concludes the paper.
%% 引言不预填实验提升数字。第四章实测完成后，可在本节最后补一句最主要发现。
%% 目前引用支撑的是具体方法关系，不宣称已穷尽相关工作或首次组合这些模块。
